\documentclass[10pt,a4paper]{article}
\pdfmapfile{+lm.map}
\pdfmapfile{+cm.map}
\pdfmapfile{+symbols.map}
\usepackage[T1]{fontenc}
\usepackage[utf8]{inputenc}
\usepackage[ngerman,provide=*]{babel}
\usepackage{lmodern,amsmath,amssymb,amsthm,booktabs,array,geometry,xcolor,hyperref,fancyhdr,microtype}
\geometry{left=23mm,right=23mm,top=24mm,bottom=24mm}
\definecolor{navy}{HTML}{17334B}
\definecolor{teal}{HTML}{087E8B}
\hypersetup{colorlinks=true,urlcolor=teal,linkcolor=navy,pdftitle={CST-DUST: Gerichteter Randvertrag und bedingte Gaussprojektion v0.1},pdfauthor={X6 Forschungsprogramm / Arbeitsausarbeitung}}
\pagestyle{fancy}\fancyhf{}
\fancyhead[L]{\small\sffamily\color{navy}CST-DUST / Randvertrag und Gaußprojektion}
\fancyhead[R]{\small\sffamily v0.1}
\fancyfoot[L]{\footnotesize 11. November 2026 / Physikalische Qualifikation: HOLD}
\fancyfoot[R]{\thepage}
\setlength{\headheight}{14pt}
\setlength{\parindent}{0pt}\setlength{\parskip}{6pt}
\setlength{\emergencystretch}{2em}
\newcommand{\R}{\mathbb R}
\newcommand{\dd}{\mathrm d}
\newcommand{\ii}{\mathrm i}
\newcommand{\sgn}{\operatorname{sgn}}
\newcommand{\sdet}{\operatorname{sqrt}_{\!c}\det}
\newcommand{\status}[1]{\textsf{\textbf{#1}}}
\newcommand{\sect}[1]{\section*{\color{navy}#1}}
\newcommand{\smallurl}[1]{{\footnotesize\url{#1}}}
\begin{document}
{\sffamily\small\color{teal}FORTSETZUNG DER MATERIEBLOCK-VORPRÜFUNG}
\vspace{5mm}

{\sffamily\bfseries\color{navy}\LARGE Gerichteter Randvertrag\par und bedingte Gaußprojektion\par}
\vspace{3mm}
{\large Ein endlicher Rechenweg mit expliziten Voraussetzungen\par}

\sect{Was dieser Schritt schließt}
Die Vorprüfung nach PR~277 identifizierte drei Hindernisse: Die vorhandenen
Geometriepolynome zählen Verklebungen, der vollständig belegte Viererrand hat
keine inneren Feldwerte, und ein Vakuum-Einteilchen-Zeuge verschwindet bei
Gesamtfeldparität. Hier werden eine gerichtete Kontrollfamilie mit einem
inneren Element und die exakte Projektion einer \emph{vorgegebenen}
quadratischen Feldwirkung auf normierte Gauß-/Hermitepakete ausgearbeitet.

Die zentrale neue Aussage betrifft singuläre Innenblöcke. Mit der
Wirkungszerlegung $H=\bigl(\begin{smallmatrix}B&J^T\\J&A\end{smallmatrix}\bigr)$
und positivem Randpaket $\Omega$ ist die vollständig projizierte,
Abel-regulierte Nullanregungsamplitude genau dann endlich im hier untersuchten
zentrierten Gaußfall, wenn
\[
\boxed{\ker A\cap\ker J^T=\{0\}.}
\]
Damit kann ein singulärer innerer Hessian trotz fehlender gewöhnlicher
Schurinverse eine endliche projizierte Amplitude besitzen. Der Beweis und eine
unabhängige Deltafunktionskontrolle stehen auf Seite~4.

\begin{center}
\begin{tabular}{p{.49\linewidth}p{.40\linewidth}}
\toprule
Ergebnis & Status / Bedeutung\\\midrule
Gerichtete Fünferelement-Kontrollfamilie & vollständig ausgezählt\\
Gaußprojektion und Nullraumkriterium & bedingt analytisch hergeleitet\\
Automatisierte Kontrollen & 27 Tests bestanden\\
Zusätzliche physikalische Materiewirkung & weiterhin offene Eingabe\\
Physikalische $8\times8$-Matrix & \texttt{NOT\_CONSTRUCTED}\\
Physikalischer Stratum-Mischblock & \texttt{NOT\_COMPUTED}\\
Alle acht physikalischen Gates & \status{HOLD}\\\bottomrule
\end{tabular}
\end{center}

\sect{Welche Entscheidungen getroffen werden}
Die gerichtete Familie ist eine \emph{ausdrücklich gewählte endliche
Kontrollfamilie}, keine Herleitung des physikalischen Historienträgers.
Gaußbreiten und Feld-Hessian werden symbolisch belassen. Die bisherigen
Uhr- und Dichteintegrale werden nicht durch unbeschränkte Feldintegrale
ersetzt. Es wird kein Übergangsterm zur Erzeugung eines gewünschten
Mischblocks hinzugefügt.

Das Ergebnis ist ein prüfbarer analytischer Anschluss, kein CST-Hamiltonian,
kein dynamischer Abschluss der acht Planungslabels und kein Bounce-Nachweis.
Die kleinen numerischen Matrizen in den Tests prüfen ausschließlich Algebra
und Integrationsformeln. Es wurden keine Produktionsläufe oder bezahlten
Rechenressourcen gestartet.

\newpage
\sect{1. Ein gerichteter Träger mit echtem Innenraum}
Es seien $L=\{0,1\}$, $I=\{2\}$, $R=\{3,4\}$. Die induzierten Ordnungen
auf $L$ und $R$ sind $X,Y\in\{A,K\}$: $A$ ist die Zweierantikette, $K$ die
fest orientierte Zweierkette. Erlaubt sind nur Relationen innerhalb der
fixierten Ränder sowie $L\to I$, $I\to R$ und $L\to R$. Die gesamte Relation
muss bereits transitiv sein; ein nachträglicher Abschluss darf die
Randordnung nicht verändern. Das innere Element muss mindestens einen
Vorgänger in $L$ und einen Nachfolger in $R$ besitzen.

Dies definiert die schwache Familie $\mathcal C^{\rm weak}_{XY}$.
Für $\mathcal C^{\rm cov}_{XY}$ kommt hinzu:
\[
\forall l\in L\ \exists r\in R:l\prec r,
\qquad \forall r\in R\ \exists l\in L:l\prec r.
\]
Jedes Randelement ist dann an die Gegenseite angeschlossen; die Verbindung
muss nicht durch Element~2 laufen. $L$ ist ein Ordnungsideal, $R$ ein Filter;
beide sind kausal konvex. Bei $AA$ in der angeschlossenen Familie sind $L$
exakt die minimalen und $R$ exakt die maximalen Elemente. Bei $K$ ist der
jeweilige Rand ein \emph{Kragen}, keine raumartige Antikette. Seine Eignung
als vollständige Randinformation bleibt zu beweisen.

\begin{center}
\begin{tabular}{lrrrr}
\toprule
$XY$ (Eingang, Ausgang) & schwach & angeschlossen & Orbits angeschl. & $\sum1/|\mathrm{Stab}|$\\\midrule
$AA$ &49&33&10&$33/4$\\
$AK$ &23&14&8&$7$\\
$KA$ &23&14&8&$7$\\
$KK$ &12&7&7&$7$\\\bottomrule
\end{tabular}
\end{center}
Die 107 schwachen bzw.~68 angeschlossenen Fälle wurden auch durch Filterung
aller 4231 gelabelten Fünferordnungen gefunden. Für $AA$ liefert eine separate
Zählung mit nichtleeren Vorgänger-/Nachfolgermengen $P,Q$:
\[
|\mathcal C^{\rm weak}_{AA}|=
\sum_{p,q=1}^2\binom2p\binom2q2^{4-pq}=49.
\]
Unter der Anschlussbedingung ergeben die Fälle $(p,q)=(1,1),(1,2),(2,1),(2,2)$
die Beiträge $20,6,6,1$, insgesamt 33. Im ersten Fall gibt es bei festem
$P,Q$ fünf zulässige Kreuzmuster, in den beiden mittleren je drei.

\textbf{Gruppenmaß.} Für austauschbare Randlabels ist
$G_{XY}=\mathrm{Aut}(X)\times\mathrm{Aut}(Y)$, mit Größen $4,2,2,1$.
Die gelabelte Summe geteilt durch $|G_{XY}|$ entspricht genau der Summe
mit Gewichten $1/|\mathrm{Stab}(C)|$. Das gilt für die Geometriesumme und
für vollständige Feldintegrale nur bei gemeinsam mittransformierten,
invarianten Randdaten und Maßen. Zusätzliche feste Marken können die Gruppe
verkleinern. Dichte-Stratumfaktoren sind hier noch nicht enthalten.

\textbf{Wirkungskontrolle.} Gespeichert sind sämtliche Ordnungs-IDs,
Stabilisatoren und die Polynome $\sum_Cz^{F_4(C)}$ mit
$F_4=N-N_0+9N_1-16N_2+8N_3$. Dies sind rohe ganzzahlige Zählsummen ohne
Staub-, Feld- und äußere Randwirkung. Sie werden nicht als neue physikalische
Übergangsmatrizen ausgegeben. Die Orientierungsumkehr ergibt eine Bijektion
$AK\leftrightarrow KA$, aber noch keine Adjungiertenrelation der Amplituden.

\newpage
\sect{2. Bedingter Wirkungs- und Zustandsvertrag}
Fixiere $C,T,w$ und die äußeren Geometriedaten. Vorausgesetzt wird eine
reell-symmetrische Feldwirkung einschließlich aller \emph{quadratischen}
Feldrandterme, ohne lineare Quelle:
\[
S_\Phi(b,q)=\frac12
\begin{pmatrix}b\\q\end{pmatrix}^{\!T}
\underbrace{\begin{pmatrix}B&J^T\\J&A\end{pmatrix}}_{H_C(T,w)}
\begin{pmatrix}b\\q\end{pmatrix},
\quad b=(u_i,u_f),\quad q\in\R^r.
\]
$u_i\in\R^{d_i}$ und $u_f\in\R^{d_f}$ sind die Randfeldkonfigurationen.
Im Fünferträger sind $d_i=d_f=2$, $r=1$. Die Formeln gelten allgemeiner.
$H$ muss aus einer vorgegebenen gemeinsamen Wirkung stammen; dieses Paket
erzeugt keine solche physikalische Wirkung. Die vorhandene Staubuhr-Hesse-Matrix
wird nicht zum zusätzlichen Materie-Hessian umbenannt.

Bei Lebesgue-Feldmaß und $\hbar>0$ wähle reelle symmetrische
$\Omega_i,\Omega_f>0$ und die \emph{normierten} Pakete
\[
\psi_{0,a}(u)=N_a\exp[-u^T\Omega_au/(2\hbar)],
\qquad N_a=\frac{(\det\Omega_a)^{1/4}}{(\pi\hbar)^{d_a/4}},
\quad a\in\{i,f\}.
\]
Die Breiten sind vorgegebene Zustandsdaten, keine Eigenwerte einer indefiniten
Raumzeit-Hesse-Matrix. Sie dürfen nicht von unbeobachteten inneren Historien
abhängig angepasst werden, wenn derselbe präparierte Randzustand gemeint ist.
Mit anderen Feldmaßen ändern sich Normierungen und Integrale entsprechend.

\textbf{Einheiten.} $H$, $\Omega$ und der nachfolgende Regulator besitzen bei
gleich skalierten Feldkoordinaten die Dimension Wirkung/Feld$^2$. Alle
Exponenten sind dimensionslos. Die Kontrollzahlen verwenden $\hbar=1$;
damit wird keine physikalische Skala festgelegt.

\sect{3. Schurformel bei invertierbarem Innenblock}
Bei invertierbarem $A$ sind
\[
E=B-J^TA^{-1}J,\qquad q_*=-A^{-1}Jb,
\qquad \Omega=\operatorname{diag}(\Omega_i,\Omega_f).
\]
Die volle reelle Innenintegration ist als Abel-/Fresnelgrenzwert definiert:
\[
\mathcal N_I=(2\pi\hbar)^{r/2}
\frac{\exp[\ii\pi\sgn(A)/4]}{\sqrt{|\det A|}}.
\]
Anschließende Gaußprojektion liefert
\[
\boxed{\mathcal A_{00}=
N_iN_f\,\mathcal N_I\,
\frac{(2\pi\hbar)^{(d_i+d_f)/2}}{\sdet(\Omega-\ii E)}.}
\]
Hier bezeichnet $\sdet$ die kontinuierlich festgelegte Wurzel auf Seite~5.
Der Faktor $\mathcal N_I$ enthält bereits die Signaturphase; sie wird nicht
nochmals als Morse-/Maslov-Faktor eingesetzt. Bei leerem Innenraum ist
$\mathcal N_I=1$ und $E=B$.

Diese Formel gilt für volle reelle Feldachsen. Ein kompaktes Innenintegral
hat nach der Verschiebung um $q_*$ ein randabhängig verschobenes Gebiet und
bedarf einer anderen Rechnung. Die kompakten Staubuhrintegrale bleiben in
jedem Fall unverändert außerhalb dieses Feldschritts.

\newpage
\sect{4. Singulärer Schurblock: das vollständige Kriterium}
Statt $A^{-1}$ einzusetzen, reguliere die inneren \emph{Feldwerte} mit
$\exp[-\epsilon q^Tq/(2\hbar)]$, $\epsilon>0$, und setze
\[
Q_\epsilon=
\begin{pmatrix}\Omega&0\\0&\epsilon I_r\end{pmatrix}-\ii H.
\]
Da der hermitesche Realteil strikt positiv ist, existiert das gemeinsame
absolut konvergente Gaußintegral. Mit $d_b=d_i+d_f$ gilt
\[
\boxed{\mathcal A_{00}(\epsilon)=
N_iN_f\frac{(2\pi\hbar)^{(d_b+r)/2}}{\sdet Q_\epsilon}.}
\]
Die Determinantenwurzel wird bei festem $\epsilon>0$ von $H=0$ stetig
fortgesetzt. Für invertierbares $A$ reproduziert der Grenzwert die
Schur-/Fresnelformel. Für singuläres $A$ kann er trotzdem existieren.

\textbf{Satz (fester, zentrierter Gaußfall).}
Unter den Voraussetzungen von Seite~3 ist
\[
Q_0\text{ invertierbar}
\quad\Longleftrightarrow\quad
\ker A\cap\ker J^T=\{0\}.
\]
Im invertierbaren Fall besitzt $\mathcal A_{00}(\epsilon)$ einen endlichen,
von null verschiedenen Grenzwert. Besitzt der gemeinsame Kern Dimension
$k>0$, enthält die Amplitude den divergenten Faktor
$(2\pi\hbar/\epsilon)^{k/2}$.

\textbf{Beweis.}
Für $z=(b,q)$ mit $Q_0z=0$ gilt
$0=\operatorname{Re}(z^\dagger Q_0z)=b^\dagger\Omega b$.
Aus $\Omega>0$ folgt $b=0$. Die beiden verbliebenen Gleichungen sind
$J^Tq=0$ und $Aq=0$. Umgekehrt ist jedes solche $q$ ein Nullvektor von
$Q_0$. Bei trivialem Kern ist $\det Q_0\ne0$; die fortgesetzte Wurzel
besitzt nahe $\epsilon=0$ einen endlichen, von null verschiedenen Grenzwert.
Bei nichttrivialem Kern entkoppeln die entsprechenden $k$ inneren Achsen
wegen der Symmetrie von $H$ von allen übrigen Variablen. Jede liefert
$\int\dd q\,e^{-\epsilon q^2/(2\hbar)}=\sqrt{2\pi\hbar/\epsilon}$.
Auf dem orthogonalen Rest ist das gemeinsame Kernkriterium wieder erfüllt.
Damit ist die Divergenzaussage gezeigt.\hfill$\square$

\textbf{Exakte Kontrolle.}
Für $b=(x,y)$, $q\in\R$, $\Omega=I_2$, $\hbar=1$ und $S_\Phi=qx$
ist $A=0$, aber $J=(1,0)$. Es gilt $\det Q_\epsilon=1+\epsilon$ und
\[
\mathcal A_{00}(\epsilon)=\frac{2\sqrt2\,\pi}{\sqrt{1+\epsilon}}
\ \longrightarrow\ 2\sqrt2\,\pi.
\]
Unabhängig liefert $\int\dd q\,e^{\ii qx}=2\pi\delta(x)$ genau denselben
projizierten Wert. Mit $J=0$ wäre das innere Nullvolumen divergent.
Das ist eine algebraische Referenz, kein gewähltes CST-Materiemodell.

\textbf{Reichweite.} Der Satz gilt punktweise bei festen $C,T,w$ und festen
positiven Paketen. Er erlaubt keinen ungeprüften Grenzwerttausch mit den
Staubintegralen. Gemeinsame Nullrichtungen sind ohne Symmetrie- und
Maßnachweis keine Eichrichtungen. Bei Quellen, nichtgaußschen Zuständen oder
kompakten Feldgebieten ist die Aussage neu zu prüfen.

\newpage
\sect{5. Phase und paritätserlaubte Randanregungen}
Für $M=\Omega-\ii E$, $\Omega>0$, $E=E^T\in\R^{d_b\times d_b}$, seien
$\lambda_j$ die Eigenwerte von $\Omega^{-1/2}E\Omega^{-1/2}$. Definiere
\[
\sdet M=(\det\Omega)^{1/2}
\exp\!\left[\frac12\sum_j\log(1-\ii\lambda_j)\right].
\]
Jeder einzelne Logarithmus ist auf der rechten Halbebene eindeutig
festgelegt. Das Produkt ist jedoch nicht allgemein die Hauptwurzel der
bereits multiplizierten Determinante.

\textbf{Vorzeichen-Gegenbeispiel.} Für $\Omega=I_4$, $E=2I_4$ ist
\[
\sdet(I_4-2\ii I_4)=(1-2\ii)^2=-3-4\ii,
\quad \sqrt{\det(I_4-2\ii I_4)}_{\rm Haupt}=3+4\ii.
\]
Eine naive Hauptwurzel würde das Vorzeichen der Amplitude umkehren.
Die gleiche Fortsetzungsregel gilt für $Q_\epsilon$, mit
$\operatorname{diag}(\Omega,\epsilon I)$ als positivem Referenzteil.

\textbf{Randwellenfunktionen.} Für einen reellen Einheitsvektor $v_a$ im
jeweiligen Randkonfigurationsraum setze
\[
\psi_{1,a}^{v_a}(u)=\sqrt{2/\hbar}\,
v_a^T\Omega_a^{1/2}u\,\psi_{0,a}(u),
\]
\[
\psi_{2,a}^{v_a}(u)=\frac{2(v_a^T\Omega_a^{1/2}u)^2/\hbar-1}{\sqrt2}
\,\psi_{0,a}(u).
\]
Dies sind normierte Hermiteanregungen der \emph{gewählten} Gaußpakete. Die
Bezeichnung physikalisches Vakuum oder Teilchen benötigt zusätzlich eine
begründete Zustands-/Modenwahl. Ein Modenprofil ist nicht selbst diese
Wellenfunktion.

Mit $a_i=(\Omega_i^{1/2}v_i,0)$,
$a_f=(0,\Omega_f^{1/2}v_f)$ und $C_b=M^{-1}$ ergeben die Gaußmomente
$\langle bb^T\rangle=\hbar C_b$ die exakten Verhältnisse
\begin{align*}
\mathcal A_{10}&=\mathcal A_{01}=0,\\
\frac{\mathcal A_{11}}{\mathcal A_{00}}&=2a_f^TC_ba_i,\\
\frac{\mathcal A_{20}}{\mathcal A_{00}}&=
\frac{2a_f^TC_ba_f-1}{\sqrt2}.
\end{align*}
Der erste Index bezeichnet den Ausgang. Für die gemeinsame regulierte
Integration ersetzt man $C_b$ durch den Randblock $(Q_\epsilon^{-1})_{bb}$.
Bei erfülltem Kriterium von Seite~4 ist auch dessen Grenzwert wohldefiniert.

Die Nullresultate folgen aus \emph{gleichzeitiger} Feldspiegelung, nicht
aus Stratum-Superselektion. Sie setzen symmetrische Wirkung, Gebiete, Maße
und geeignete Präparation voraus. Die geraden Kandidaten $00$ und $11$
sind erlaubt, aber nach Summe über Historien und Staubintegrale nicht
garantiert ungleich null. Die Erweiterung um $20$ vergrößert ausdrücklich
den Kontrollraum über die bisherigen zwei Materielabels hinaus.

\newpage
\sect{6. Stratumwechsel benötigt definierte Atomdaten}
Der Unterschied zwischen einem Dichte-Nullatom und einem Materievakuum
bleibt wesentlich. Für eine einzelne Dichtekoordinate ist das geerbte Maß
\[
\dd\nu(w)=\tfrac12\delta_0(\dd w)+\tfrac12\mathbf1_{(0,1)}(w)\dd w.
\]
Die normierten konstanten Stratumfunktionen sind
$e_0=\sqrt2\,\mathbf1_{\{0\}}$ und
$e_+=\sqrt2\,\mathbf1_{(0,1)}$. Für einen bezüglich $\nu$ definierten
Integralkern $K(w_f,w_i)$ folgt daher
\[
\langle e_0,Ke_+\rangle_\nu=
\frac12\int_0^1 K(0,w_i)\dd w_i,
\qquad
\langle e_+,Ke_0\rangle_\nu=
\frac12\int_0^1 K(w_f,0)\dd w_f.
\]
Die beiden Ausdrücke brauchen nicht gleich oder komplex konjugiert zu
sein. Dafür wäre eine Adjungierten- bzw.~Orientierungsrelation des
\emph{vollständigen} Kerns nachzuweisen.

\textbf{Bedingte Unterbestimmtheit.}
Kennt man $K$ nur fast überall auf $(0,1)^2$ bezüglich Lebesgue-Maß,
sind die Atomzeile und Atomspalte nicht festgelegt. Die Ergänzung
\[
K_c=K+c\,\mathbf1_{\{0\}}(w_f)\mathbf1_{(0,1)}(w_i)
+\overline c\,\mathbf1_{(0,1)}(w_f)\mathbf1_{\{0\}}(w_i)
\]
ändert den positiven Bulk nicht, aber die normierten Mischblöcke um
$c/2$ bzw.~$\overline c/2$. Sie bewahrt die Hermitizität, falls diese
vorher vorlag. Bulkwerte und Hermitizität bestimmen diese Kopplung also
nicht allein. Stetigkeit bis zum Atom oder eine auf sämtlichen Strata
definierte gemeinsame Wirkung könnten die Mehrdeutigkeit beseitigen;
diese zusätzlichen Voraussetzungen müssen ausdrücklich vorliegen.

Dies ist \emph{kein Vorschlag}, $c$ frei einzusetzen. Es erklärt, warum
eine tatsächliche atomare Randamplitude aus der vollständigen
Wirkungs-/Maßdefinition gewonnen werden muss. Die Hilfsrelation $wf=0$
am Atom beseitigt nur Terme mit genau diesem Faktor. Kopplungen an
positive Gewichte anderer Elemente verschwinden dadurch nicht automatisch.

\textbf{Mehrere Randdichten.} Im Zweierkragen gibt es bereits die vier
Masken $00,0+,+0,++$. Das Label ``$+$'' ist ohne Auswahl oder Summe dieser
Masken nicht eindeutig. Es wird hier nicht stillschweigend mit ``beide
positiv'' gleichgesetzt. Die acht bisherigen Planungslabels bleiben in
ihrer Reihenfolge erhalten, sind aber noch keine aus dieser neuen
Randfamilie abgeleitete physikalische Basis.

\textbf{Was ein gemischter Block schließlich enthalten muss.}
Unter einem vollständig festgelegten Vertrag lautet der Zieltyp
\[
B_{\beta\alpha}(\epsilon)=\frac1{|G_{\rm fix}|}
\sum_C\int\dd\mu_d(T,w)\,
\overline{f_\beta(T_R,w_R)}f_\alpha(T_L,w_L)
e^{\ii(S_g+S_d+S_{\partial,0})/\hbar}
\mathcal A_{n_fn_i}(C,T,w;\epsilon).
\]
Quadratische Feldrandterme stecken bereits in $H$; $S_{\partial,0}$ umfasst
hier nur feldunabhängige äußere Randterme. Diese Zerlegung verhindert
Doppelzählung. Der globale $\epsilon\downarrow0$-Grenzwert, Gramstruktur
und Trunkierungsfehler bleiben gesondert zu kontrollieren.

\newpage
\sect{7. Prüfung, offene Eingaben und nächster Engpass}
Die neue Rechnung schließt die gerichtete \emph{Kontrollzählung}, die
bedingte Gaußprojektion, die Phasenwahl und das gemeinsame Nullraumkriterium.
Die 27 automatisierten Tests decken unter anderem ab:
\begin{itemize}
\setlength\itemsep{1pt}
\item vollständigen Abgleich mit allen 4231 gelabelten Fünferordnungen,
      Randextrema, Gruppenorbits, Wirkungsinvarianz und Orientierungsumkehr;
\item exakte stationäre Schuridentität, Quellenableitungen der Gaußmomente
      und neun Orthogonalitäts-/Normierungsintegrale der Hermitepakete;
\item unabhängige Gauß-Hermite-Quadratur im Rand- und vollständigen
      regulierten Feldraum mit absoluter Abweichung unter $10^{-11}$;
\item singulär gekoppelten endlichen Fall, entkoppeltes divergentes
      Nullvolumen und das Vorzeichen-Gegenbeispiel der Determinantenwurzel;
\item Eingabenkontrollen, Basispermutation, Paritätsblöcke und einen
      Gegenfall zum Abschluss der Besetzungen $0,1$; physikalische Werte bleiben offen.
\end{itemize}
Die Quadraturkontrollen sind numerische Implementierungstests, keine
zertifizierten Fehlerintervalle einer physikalischen Amplitude. Die
allgemeinen Aussagen werden durch die Beweise im Text getragen.

\textbf{Reproduktion.} Im entpackten Verzeichnis ausführen:
\begin{verbatim}
python analyze.py
python -m unittest discover -s . -p 'test_*.py' -v
python verify.py
\end{verbatim}
\texttt{verify.py} prüft die gespeicherten SHA-256-Hashes und reproduziert
beide JSON-Ergebnisdateien in einem frischen Verzeichnis bytegleich.
\texttt{gaussian.py} enthält nur bedingte Projektionsroutinen und keine
Erzeugung physikalischer CST-Materieblöcke. \texttt{physical\_inputs.json}
hält die fehlenden Daten als \texttt{null}; sie werden nicht mit Nullen
oder Testwerten gefüllt.

\textbf{Nächster wissenschaftlicher Engpass.}
Jetzt wird eine unabhängig begründete zusätzliche Feldwirkung samt
Randfunktional und Zustandswahl benötigt. Sie muss auf \emph{derselben}
markierten Geometrie und auf allen benutzten Dichtestrata definiert sein.
Erst daraus werden $H_C(T,w)$, Randbreiten bzw.~komplexe Struktur und
die Masken des Labels ``$+$'' gewonnen. Danach folgt ein gemeinsamer,
regulierter gerader Mischblock mit Fehlerbudget und identischem
Entkopplungsvergleich. Ein frei ergänztes $q_{0+}$ wäre keine Lösung
dieses Herleitungsproblems.

\sect{Herkunft und Literaturabgleich}
Projektbasis: PR~277 wurde auf \texttt{main} zusammengeführt. Gepinnter Anschluss:
\par{\footnotesize\texttt{fdf1deb6ac071056d533a0937e21c238f42b2974}}.
\par Der unveränderte Enumerationsbaustein \texttt{gluing.py} wurde gegen diesen Git-Blob abgeglichen:
\par{\footnotesize\texttt{d4c9ebdf640040a5fa1f00042d709c7b572aac5e}}.
\par Die lokale Vorprüfung ist als \path{inputs/SOURCE_PREFLIGHT.md} beigefügt.

{\small
\begin{enumerate}
\setlength\itemsep{1pt}
\item S. Aslanbeigi, M. Saravani, R. D. Sorkin, \emph{Generalized Causal Set
d'Alembertians} (2014). Retardierte, nichtlokale Operatoren; kein Ersatz
für die hier fehlende symmetrische Feldwirkungsdefinition.
\smallurl{https://arxiv.org/abs/1403.1622}
\item R. D. Sorkin, \emph{Scalar Field Theory on a Causal Set in Histories Form}
(2011/2017); Nomaan X, \emph{Quantum Field Theory On Causal Sets} (2023).
Historien- und Zustandskonstruktionen als getrennte physikalische Eingaben.
\smallurl{https://arxiv.org/abs/1107.0698}\quad
\smallurl{https://arxiv.org/abs/2306.04800}
\item M. Buck et al., \emph{Boundary Terms for Causal Sets} (2015);
F. Dowker, R. Liu, D. Lloyd-Jones, \emph{Timelike boundary and corner terms
in the causal set action} (2025). Unterschiedliche Randprobleme; daraus
folgt kein bereits qualifiziertes Randgesetz für diesen Kontrollträger.
\smallurl{https://arxiv.org/abs/1502.05388}\quad
\smallurl{https://arxiv.org/abs/2501.00139}
\end{enumerate}}
\newpage
\sect{8. Abgleich der nachgereichten Formeln}
Die zusätzlichen Bilder 081--099 wurden in dieser Fortsetzung geprüft.
Sie enthalten $L_C=B_C+m_\Phi^2W_C$, die Rand-/Innenzerlegung,
das verallgemeinerte Eigenwertproblem und den Ausdruck
$M_C^{nm}=e_n^\dagger L^{\rm eff}_{\partial,C}e_m$.
Ihre Einordnung ist wie folgt.

\textbf{1. Operator und Wirkungs-Hessian nicht vermischen.}
Ist $L_C$ bereits die reell-symmetrische Matrix der Bilinearform
$S=\frac12\Phi^TL_C\Phi$, kann sie als $H$ in diesem Bericht dienen.
Ist $B_C$ dagegen ein retardierter, nicht symmetrischer Operator, wird
$B_C+m_\Phi^2W_C$ durch den Massenterm nicht automatisch symmetrisch.
Es gilt
\[
\nabla_\Phi\bigl(\tfrac12\Phi^TL_C\Phi\bigr)
=\tfrac12(L_C+L_C^T)\Phi.
\]
Soll $L_C$ ein bezüglich $W_C$ selbstadjungierter \emph{Operator} sein,
wäre etwa $H=W_CL_C$ die zugehörige Bilinearform. Das ist eine andere
Konvention. Vorzeichen, Einheiten und Randterme bleiben zu definieren;
eine Symmetrisierung wird hier nicht als physikalische Herleitung ausgegeben.

\textbf{2. Die Schurformeln sind bedingt korrekt.}
$\Phi_I^*=-L_{II}^{-1}L_{I\partial}\Phi_\partial$ und
$L^{\rm eff}_\partial=L_{\partial\partial}
-L_{\partial I}L_{II}^{-1}L_{I\partial}$ gelten für die stationäre
Elimination eines invertierbaren, korrekt definierten Wirkungsblocks.
Das vollständige Quantenintegral braucht zusätzlich Phase, Determinante,
Maß und Integrationsvorschrift. Seite~4 behandelt die singuläre Erweiterung.

\textbf{3. Im eigenen Eigenprofilraum ist die Projektion diagonal.}
Falls $L^{\rm eff}_\partial e_m=\lambda_mW_\partial e_m$ und
$e_n^\dagger W_\partial e_m=\delta_{nm}$, folgt unmittelbar
\[
\boxed{M_C^{nm}=\lambda_m\delta_{nm}.}
\]
Das Verschwinden der Nebendiagonalen ist hier eine Folge der Basiswahl,
kein Nachweis fehlender quantisierter Materieübergänge. In einer festen
Referenzprofilbasis können Nebendiagonalen auftreten; auch diese sind
zunächst Koeffizienten einer Feld-Bilinearform. Sie werden erst durch
definierte Wellenfunktionen und die gemeinsame Amplitude zu den auf
Seite~5 berechneten Übergangsmatrixelementen. Modenindex $k$ und
Besetzungszahl $n$ müssen verschieden bezeichnet werden.

Falls $W_\partial$ direkt aus Dichtegewichten gebildet wird, kann es am
Nullatom singulär werden. Dann ist eine vollständige Normierung
$e_n^\dagger W_\partial e_m=\delta_{nm}$ auf diesem Feldraum unmöglich.
Ein Quotient oder ein anderes positives Randprodukt braucht eine eigene
Begründung. $W_\partial$ als Profilmetrik ist zudem nicht automatisch das
Integrationsmaß für Randwellenfunktionen.

\textbf{4. Die Summe der stationären Wirkungen ist noch keine volle
Quanten-Randwirkung.}
Schreibt man den eingesandten Ausdruck als
$\mathcal W_{C,s}=S_{\rm BDG}+S_d+\frac12b^TEb$, so fehlt im unbeschränkten,
invertierbaren Gaußfall gegenüber $-\ii\hbar\log K_C(b)$ insbesondere
\[
-\ii\hbar\log\mathcal N_I(C,T,w).
\]
Eine Logarithmusbranche ist festzulegen. Der Zusatz kann von Geometrie
und Staubdaten abhängen und ist daher innerhalb der gemeinsamen Summe
nicht ohne Nachweis eine irrelevante Konstante. Bei kompakten Feldgebieten
kommen randabhängige Integralfaktoren hinzu. Für den vollständigen
Historienkern ist außerdem der Logarithmus \emph{nach} der gemeinsamen
Summe und Integration zu bilden, sofern der Kern dort nicht null ist.
\newpage
\sect{9. Acht Zustände: Indizes und Blockanordnung}
Die ergänzten Bilder 071--080 geben die Labels $|G,s,n\rangle$ und eine
$8\times8$-Anordnung vor. Diese Größe passt zu zwei Geometrien, zwei
\emph{ausgewählten} Stratumzuständen und zwei \emph{ausgewählten}
Materiewellenfunktionen. Sie ist keine Dimensionsaussage über den
vollständigen physikalischen Zustandsraum. Insbesondere besitzt schon
$L^2((0,1),\dd w)$ mehr Zustände als das konstante Paket $e_+$.

Für die Blockmatrix
\[
K^{\rm geom}=\begin{pmatrix}K_{AA}&K_{AK}\\K_{KA}&K_{KK}\end{pmatrix}
\]
gilt hier ausdrücklich: \textbf{Zeile = Ausgang, Spalte = Eingang}.
Innerhalb jedes Geometrieblocks lautet die Reihenfolge
$(s,n)=(0,0),(0,1),(+,0),(+,1)$. Damit ist das eingesandte Schema
\[
K_{XY}=\left(\begin{array}{cccc}
K_{00}^{00}&K_{00}^{01}&K_{0+}^{00}&K_{0+}^{01}\\
K_{00}^{10}&K_{00}^{11}&K_{0+}^{10}&K_{0+}^{11}\\
K_{+0}^{00}&K_{+0}^{01}&K_{++}^{00}&K_{++}^{01}\\
K_{+0}^{10}&K_{+0}^{11}&K_{++}^{10}&K_{++}^{11}
\end{array}\right)_{XY}
\]
ein konsistentes \emph{Indexschema}. Beide Indexpaare laufen jeweils als
Ausgang--Eingang. Unbekannte Einträge bleiben unbekannt.

\begin{center}
\begin{tabular}{cll}
\toprule
Position & Geometrie zuerst (Bilder) & bisheriger Planungsvertrag\\\midrule
1 &$|A,0,0\rangle$ &$|A,0,0\rangle$\\
2 &$|A,0,1\rangle$ &$|A,0,1\rangle$\\
3 &$|A,+,0\rangle$ &$|K,0,0\rangle$\\
4 &$|A,+,1\rangle$ &$|K,0,1\rangle$\\
5 &$|K,0,0\rangle$ &$|A,+,0\rangle$\\
6 &$|K,0,1\rangle$ &$|A,+,1\rangle$\\
7 &$|K,+,0\rangle$ &$|K,+,0\rangle$\\
8 &$|K,+,1\rangle$ &$|K,+,1\rangle$\\\bottomrule
\end{tabular}
\end{center}
Mit $p=(1,2,5,6,3,4,7,8)$ und $P_{r,p_r}=1$ gilt
\[
K^{\rm plan}=P K^{\rm geom}P^T.
\]
Gram-Matrizen und weitere Zustandsoperatoren sind ebenso umzuschreiben.
Das ändert keine Amplitude und überschreibt die bisherige Reihenfolge nicht.

\textbf{Orientierung der Träger.} Abschnitt~1 und die Zähldatei verwenden
$XY$ als \emph{Eingang--Ausgang}. Somit gehört ein Matrixblock $K_{XY}$ zur
dortigen Trägerfamilie $\mathcal C_{YX}$. Die gleichen Zählwerte für $AK$
und $KA$ würden eine Verwechslung verdecken. Sie liefern insbesondere
keinen Beweis für $K_{AK}=K_{KA}^\dagger$.

\textbf{Moden sind keine Besetzungen.} Die Größe acht setzt hier genau
zwei Wellenfunktionen pro $(G,s)$ voraus, etwa $n=0,1$ für eine gewählte
Anregungsrichtung. Zwei unabhängige Moden mit jeweils Besetzungen $0,1$
ergäben vier Materiezustände und insgesamt 16 Labels. Beschränkt man
diese beiden Moden dagegen auf höchstens eine Anregung insgesamt,
sind es drei Materiezustände und 12 Labels.

\newpage
\sect{10. Projektion, Parität und fehlender dynamischer Abschluss}
Unter der Gesamtfeldparität von Abschnitt~5 hat das Paket $n$ die Parität
$(-1)^n$. Deshalb verschwinden die Einträge $K_{ab}^{01}$ und
$K_{ab}^{10}$ für alle Geometrien und Strata, \emph{falls} Wirkung, Maß,
Integrationsgebiet und Regulator diese Symmetrie besitzen. Sortiert nach
Parität zerfällt das Achtzustandsschema dann in zwei $4\times4$-Blöcke.
Stratumwechsel innerhalb eines solchen Blocks bleiben erlaubt.

Die Gleichsetzung $K_{XY}=K^{\rm grav}_{XY}I_4$ wäre eine wesentlich stärkere
Zusatzannahme: keine Stratum- oder Materieübergänge und derselbe Faktor für
alle vier ausgewählten Zustände. Gesamtfeldparität begründet diese
Faktorisierung nicht. Auch eine rein quadratische Feldwirkung garantiert
keine Erhaltung der Besetzungszahl einer gewählten Referenzbasis.
Der auf Seite~5 bereits berechnete Ausdruck $\mathcal A_{20}$ kann ungleich
null sein. \textbf{Parität schließt den Raum $n=0,1$ daher nicht ab.}

\textbf{Präziser bedingter Projektionsvertrag.} Falls ein physikalischer
Hilbertraum und ein Übergangsoperator $U:\mathcal H_i\to\mathcal H_f$
begründet sind, seien $V_a:\mathbb C^8\to\mathcal H_a$ isometrische
Einbettungen der ausgewählten Zustände und $P_a=V_aV_a^\dagger$.
Die projizierte Matrix und der ausgehende Rest sind
\[
K_8=V_f^\dagger U V_i,
\qquad R=(I-P_f)UV_i.
\]
Ein dynamischer Abschluss verlangt $R=0$ bzw.~eine begründete kleine
Fehlerschranke. Unter der \emph{zusätzlichen} Annahme $U^\dagger U=I$ gilt
\[
\boxed{I_8-K_8^\dagger K_8=R^\dagger R,}
\qquad
1-\|K_8x\|^2=\|Rx\|^2\quad(\|x\|=1).
\]
Dies folgt direkt aus $P_f+(I-P_f)=I$. Der Satz ist ein Prüfwerkzeug,
keine Behauptung einer bereits vorhandenen unitären CST-Zeitentwicklung.
Bei einer Constraint-/Projektoramplitude ist stattdessen deren eigenes
Skalarprodukt und Verklebungsgesetz zu prüfen.

Sind die ausgewählten Zustände zunächst nicht orthonormal, bezeichnet
$B_{\beta\alpha}=\langle v_{f,\beta},Uv_{i,\alpha}\rangle$ nur ihre
Matrixelemente. Bei positiven Gram-Matrizen $G_i,G_f$ lautet die orthonormale
Darstellung $\widehat B=G_f^{-1/2}BG_i^{-1/2}$. Eine singuläre Gram-Matrix
benötigt zuerst eine begründete Reduktion; bloßes Umsortieren genügt nicht.

\textbf{Vervollständigung des eingesandten Rechenwegs.}
Für invertierbaren Innenblock ist der bedingte Weg
\[
(C,T,w;\text{Wirkung und Maß})\longrightarrow H_C
\longrightarrow(E_C,\mathcal N_{I,C})
\longrightarrow\mathcal A_C^{n_fn_i}
\longrightarrow K_8.
\]
Die Randwellenfunktionen einschließlich Breiten und Modenrichtung sind vor
der Projektion zusätzliche Eingaben. Der letzte Pfeil enthält die gemeinsame
Historien-/Staubsumme, Atomdaten und Normierung aus Abschnitt~6.
Für singuläre Innenblöcke ersetzt $Q_\epsilon$ den gewöhnlichen Schurschritt.
Ein Schurmatrix-Eigenprofil allein liefert noch keinen quantisierten
Randzustand; der Determinantenfaktor darf nicht verloren gehen.

\textbf{Prüfnachweis.} Drei zusätzliche exakte Kontrollen prüfen alle
64 Indexzuordnungen, die Paritätsblockstruktur und einen paritätserhaltenden
unitären Dreizustands-Gegenfall mit Übergang $0\to2$.
\path{basis_contract.json} dokumentiert Konventionen und offene Eingaben.
Alle acht physikalischen Gates bleiben \status{HOLD}.
\end{document}
